\documentclass[11pt,a4paper]{article}
\usepackage[utf8]{inputenc}
\usepackage[T1]{fontenc}
\usepackage[spanish,es-nodecimaldot]{babel}
\usepackage{lmodern}
\usepackage[margin=2.5cm]{geometry}
\usepackage{microtype}
\usepackage{csquotes}
\usepackage[backend=biber,style=ieee,sorting=none]{biblatex}
\addbibresource{referencias.bib}
\usepackage[hidelinks]{hyperref}
\usepackage{xurl}
\setlength{\parindent}{0pt}
\setlength{\parskip}{0.55em}
\hypersetup{pdftitle={Amazon.com: del monolito a una arquitectura orientada a servicios}}
\begin{document}
\begin{center}
{\small Aplicaciones Distribuidas -- FORM-IC-01\par}
\vspace{0.4em}
{\Large\bfseries Amazon.com: del monolito a una arquitectura orientada a servicios\par}
\vspace{0.4em}
{\small Autor: Nixon Piguave\par}
\end{center}

\section*{Delimitación del caso}
La investigación estudia la evolución de la plataforma de comercio electrónico Amazon.com hacia una arquitectura orientada a servicios (SOA). El foco es la separación de responsabilidades entre servicios y sus equipos. Se examina la transformación de la tienda; AWS aporta documentación técnica sobre ese proceso.

\section*{Descripción y escala}
La carta a los accionistas correspondiente a 2010 identifica explícitamente SOA como base de las tecnologías de Amazon. Describe servicios que colaboran en recomendaciones, gestión de pedidos e inventario. Para construir una página de detalle de producto se invocaban entre 200 y 300 servicios. Esta cifra corresponde al periodo descrito en la carta y permite dimensionar el caso sin asumir que representa la arquitectura actual~\cite{amazon2010letter}.

\section*{Arquitectura anterior y transformación}
El apéndice B de la guía técnica de AWS sitúa en 2001 una aplicación monolítica que reunía presentación, lógica de negocio y acceso a datos. La transformación comprendió cambios tecnológicos, organizativos y de operación: equipos responsables de productos y herramientas compartidas para automatizar entregas. La guía también explica que un estudio interno de 2009 detectó demoras debidas a procesos manuales, lo que impulsó el desarrollo de herramientas de canalización de código~\cite[apéndice B, pp. 17--19]{aws2021cee}.

\section*{Beneficios y desafíos identificados}
La separación mediante interfaces permitió que los servicios evolucionaran con menor impacto sobre otros componentes~\cite{amazon2010letter}. La guía de AWS informa de 50 millones de despliegues en 2016 como resultado de la transformación conjunta de herramientas, procesos y organización. Esa cifra no demuestra por sí sola el efecto aislado de SOA. Entre los aspectos operativos abordados estuvieron el diseño ante fallos, la supervisión centralizada y la responsabilidad de cada equipo sobre la construcción y operación de su producto~\cite[apéndice B, pp. 18--19]{aws2021cee}.

\clearpage
\printbibliography[title={Referencias}]
\end{document}